\chapter{Web and DOM Renderers}
\label{ch:web-dom-renderers}
\label{ch:canvas-ascii-backends}

Four CNA identities render XNA-style sprites in a browser without exposing CNA's WebGL renderer
contract. \texttt{CANVAS}
paints pixels into a Canvas2D context.  \texttt{HTML\_DOM} represents sprites as pooled,
CSS-transformed elements. \texttt{SVG\_DOM} emits SVG image nodes and filters. \texttt{PIXIJS}
uses the pinned PixiJS~7 UMD runtime over the browser's own rendering choice. All are
Emscripten-only and 2D-only, but their browser object models and evidence differ sharply.

\begin{longtable}{p{0.18\textwidth}p{0.32\textwidth}p{0.20\textwidth}p{0.17\textwidth}}
\toprule
Identity & Browser representation & Additive blend & RT / MRT / query \\
\midrule
\endfirsthead
\toprule
Identity & Browser representation & Additive blend & RT / MRT / query \\
\midrule
\endhead
\texttt{CANVAS} & \texttt{CanvasRenderingContext2D} pixels & reported true & yes / no / no \\
\texttt{HTML\_DOM} & pooled \texttt{div} elements with CSS transforms & browser-dependent \texttt{plus-lighter} & yes / no / no \\
\texttt{SVG\_DOM} & SVG \texttt{image} nodes and \texttt{feColorMatrix} & browser-dependent \texttt{plus-lighter} & yes / no / no \\
\texttt{PIXIJS} & PixiJS sprites and render textures & Pixi blend mode & yes / no / no \\
\bottomrule
\end{longtable}

\section{CANVAS: immediate raster operations}

The renderer obtains \texttt{canvas.getContext('2d')} and maps texture uploads and sprite
operations onto \texttt{drawImage}, \texttt{putImageData}, and \texttt{fillRect}.  A 2D render
target is another Canvas2D surface.  There is no shader stage, depth buffer, MRT, occlusion
query, or 3D path; inherited effect-aware draws terminate in the explicit 3D refusal.

Only additive blending is reported among the 13 capability flags.  That narrow switch is more
useful than the renderer contract's default-true policy because callers do not have to discover
the 2D ceiling by constructing unsupported objects.

\subsection{Built tests are not registered tests}

Canvas test executables exist, but the CMake files deliberately do not register them with
CTest.  A green host test run therefore says nothing about Canvas pixels.  The meaningful path
is an Emscripten build executed in a browser with observable Canvas output.

Node alone is not a substitute: it supplies neither a browser canvas nor the rendering event
loop expected by the module.  Headless Chromium can provide those APIs, but only if the test
actually opens the page, drives frames, and reads pixels or a screenshot.

\section{HTML\_DOM: sprites as pooled CSS objects}

\texttt{HTML\_DOM} does not rasterize every sprite into one surface.  It maintains a pool of
DOM elements, assigns image content, and expresses destination position, scale, rotation,
origin, opacity, crop, and visibility through CSS and element properties.  Reusing nodes avoids
unbounded allocation as the visible sprite set changes frame to frame.

This representation makes browser layout and compositing part of the renderer contract.
Ordering becomes DOM stacking order, transforms use CSS conventions, and additive blending
depends on browser support for \texttt{plus-lighter}.  The renderer therefore cannot promise
one identical native blend implementation across browser versions.

Its off-screen target is DOM-backed, not a GPU framebuffer.  MRT and queries are absent.  The
inherited 3D route reaches a local unconditional throw and bypasses the common
\texttt{WarnAndStub} policy, so diagnostic configuration cannot turn that refusal into a
warning-only placeholder.

\subsection{The strongest browser evidence in this cluster}

HTML DOM has a host contract target and a real Playwright/Chromium CI path under Xvfb.  This is
stronger than merely compiling Emscripten output because it executes the browser object model.
It still matters what the script observes: node count and style strings prove a different layer
from final composited pixels.

\section{SVG\_DOM: retained vector structure}

\texttt{SVG\_DOM} creates SVG image nodes and uses SVG transforms for sprite placement.
Tinting is expressed through a real \texttt{feColorMatrix} filter rather than by rewriting
texture pixels. Additive sprites set \texttt{mix-blend-mode: plus-lighter}; the capability query
returns the running browser's \texttt{CSS.supports} answer, just as HTML\_DOM does. The retained
tree is inspectable and naturally matches vector composition, but it inherits SVG's filter,
sampling, blending, and browser-layout semantics.

Like HTML DOM, it owns a DOM off-screen target, has no MRT or queries, and its 3D refusal
bypasses \texttt{WarnAndStub}. Its browser smoke page checks raw CSS support, the public
capability answer and the emitted sprite style together, and separate screenshot fixtures cover
pixels and scissor order. Its verification tier nevertheless remains weaker: the native host
target is behind an off-by-default \texttt{CNA\_BUILD\_SVG\_DOM\_HOST\_TESTS} option, and no CI
workflow drives either it or the browser pages. A runnable browser script is a real test route,
not evidence that the route currently gates changes.

\section{PIXIJS: browser library objects and render textures}

\texttt{PIXIJS} is Emscripten-only and resolves a pinned PixiJS~7.4.2 UMD artifact with a recorded
SHA-256; automatic download is on unless \texttt{CNA\_PIXIJS\_ROOT} supplies the file. It maps
textures, SpriteBatch/SpriteFont draws, transforms, scissor, blend state and one active
\cnaclass{RenderTarget2D} to PixiJS objects. The browser library may use WebGL internally, but the
CNA identity remains PIXIJS and its contract is the 2D adapter, not CNA's WEBGL2 family.

MRT, mip levels above zero, 3D, custom/compiled effects, cube and volume textures, cube targets
and occlusion queries are explicit refusals. Render textures make its off-screen route stronger
than NanoVG's, while more than one simultaneous target still throws. Browser execution must prove
the resolved library, object lifetime and final pixels; a successful Emscripten link alone does
not.

\section{The removed ASCII identity}

Older editions paired Canvas with an \texttt{ASCII} renderer.  That public identity and its
selector are gone.  The reusable quantization logic now lives in
\cnaclass{CNA::Graphics::AsciiPostProcessEffect}; Appendix~\ref{ch:cnaext-catalog} records the
current CNAEXT surface.  Historical architecture is not a fourth row in the live renderer
registry.

\section{Comparison by object model}

The four families can display an identical sprite scene while failing in different places:

\begin{itemize}
  \item Canvas failures cluster around pixel upload, context state, clipping, and immediate
        compositing;
  \item HTML DOM failures cluster around pooling, CSS transforms, stacking, and browser blend
        support;
  \item SVG DOM failures cluster around retained-node lifetime, SVG transforms, filters, and
        sampling.
  \item PixiJS failures cluster around UMD artifact resolution, Pixi object/render-texture
        lifetime, batching, and the browser driver selected below the library.
\end{itemize}

A useful cross-renderer scene therefore includes crop, origin-centred rotation, alpha and
additive overlap, clipping, target bind/unbind, and node-count stability over several frames.
Inspect both structure and final pixels: the DOM can look mechanically correct while the
browser composites a different color, and a correct screenshot can hide unbounded node growth.

\section{Deployment checklist}

Record the CNA identity, Emscripten version, browser and browser version, device-pixel ratio,
canvas/SVG dimensions, and whether execution was visible or headless.  Confirm that the page
reaches multiple animation frames and that the observable is meaningful for the feature under
test.  Browser compilation is only the first gate; these renderers exist in the DOM that runs
after it.
